% ============================================================================
% AION Functional Glyph Memory and Mastery Reconstruction
% ============================================================================
\section{Functional Glyph Memory and Source-Closed Mastery Reconstruction}
\label{sec:functional-glyph-memory}

This stage addressed a distinction that had previously remained hidden by
aggregate competency evidence: remembering that a method exists is not the
same as retaining enough verified structure to reconstruct and apply that
method on unfamiliar work.  The programme therefore moved from shallow topic
glyphs and accumulated assessment receipts toward a typed, provenance-bound
functional memory whose contents must survive source-closed reconstruction,
fresh execution, counterexamples, transfer, and safety testing.

The development sequence comprised five linked steps:
\begin{enumerate}
    \item invention of a general sealed-artifact verifier atom;
    \item a fresh Advanced/Expert reconstruction audit that failed closed at
          five of six artifacts;
    \item construction of the Functional Glyph Lexicon and its typed method
          compiler;
    \item successful repetition of the unchanged failed examination; and
    \item autonomous induction of new functional procedures from verified
          competency trajectories.
\end{enumerate}

The resulting memory is neither a raw answer store nor a conventional vector
index.  A glyph acts as a compact address into a typed graph containing
meaning, procedure, constraints, failures, verification, transfer, and
immutable evidence.

\subsection{Recursive Sealed-Artifact Verification}

Before a functional memory could be trusted, AION required an authority able
to test heterogeneous reconstructed artifacts without granting them ambient
execution rights.  AION invented and retained the verifier atom
\begin{verbatim}
EXECUTE_SEALED_ARTIFACT_CONTRACT
\end{verbatim}
under the promoted procedure
\begin{verbatim}
procedure_recursive_verifier_atom_invention_v1
\end{verbatim}

The atom independently verifies three artifact families:
\begin{itemize}
    \item structured natural-language artifacts;
    \item recoverable dependency graphs; and
    \item Python programs executed in disposable isolated processes.
\end{itemize}

It transferred across all three families, rejected all six adversarial
counterexamples, and received no network, shell, credential, or live-write
authority.  No candidate was installed before the complete verification gate
passed.  The focused verifier stack passed nine of nine tests.

\begin{table}[htbp]
\centering
\caption{Recursive verifier-atom results.}
\label{tab:recursive-verifier-atom}
\begin{tabular}{lr}
\hline
\textbf{Measure} & \textbf{Result} \\
\hline
Artifact-family transfer & $3/3$ \\
Adversarial counterexamples rejected & $6/6$ \\
Network, shell, or live-write authority & $0$ \\
Premature installations & $0$ \\
Focused verification tests & $9/9$ \\
\hline
\end{tabular}
\end{table}

\subsection{The First Source-Closed Competency Audit}

AION was then evaluated across five subjects that the progressive ledger
described as Advanced or Expert: English, Software Engineering, Algorithms and
Data Structures, Python, and Testing and Debugging.  No previous solution was
provided, no learning source was reopened, and the proposal substrate received
only compressed verified memory.  It generated six new artifacts.

\begin{table}[htbp]
\centering
\caption{Initial source-closed reconstruction audit.}
\label{tab:initial-reconstruction-audit}
\begin{tabular}{p{0.68\linewidth}c}
\hline
\textbf{Fresh artifact} & \textbf{Result} \\
\hline
Letter to the creator, 154 words & Pass \\
Software-engineering explanation, 114 words & Pass \\
Software-engineering essay, 530 words & Pass \\
Technical knowledge reconstruction & $4/4$ \\
Persistent multi-goal plan & Pass \\
Executable priority planner and restartable mission runner & Fail \\
\hline
\end{tabular}
\end{table}

The planning control plane itself produced 20 mission nodes, five milestones,
four learning dependencies, 11 measurable work stages, recovery coverage for
every work stage, and an acyclic dependency graph.  The failure occurred in
fresh practical construction: neither tested local proposal model reliably
reconstructed both required Python interfaces in one complete executable
module.  The procedure
\begin{verbatim}
procedure_open_competency_reconstruction_mission_v1
\end{verbatim}
therefore remained \texttt{REJECTED\_OR\_INCOMPLETE} at $5/6$.

This negative result exposed the central memory defect.  Existing glyphs
preserved concepts, evidence pointers, and coverage more successfully than
they preserved executable procedural knowledge.  AION could recall notions
such as graphs, topological ordering, rollback, or testing, but could not
always reconstruct the input schema, state transition rules, invariants,
failure cases, and verification properties needed to build a working solution.
The test was not weakened and the Expert label was not allowed to substitute
for fresh performance.

\subsection{Functional Glyph Lexicon}

The repair was a layered Functional Glyph Lexicon shared across language,
programming, science, planning, and action.  It separates eight node types:
\begin{enumerate}
    \item lexical senses and aliases;
    \item concepts and entities;
    \item reconstructable procedures;
    \item invariants and constraints;
    \item failure signatures and prior mistakes;
    \item executable verifiers;
    \item source-disjoint transfer abstractions; and
    \item immutable provenance records.
\end{enumerate}

The graph uses typed relationships including
\texttt{requires}, \texttt{uses}, \texttt{must\_preserve},
\texttt{fails\_as}, \texttt{verified\_by}, \texttt{transfers\_to},
\texttt{supported\_by}, \texttt{composes\_with},
\texttt{language\_hosts}, \texttt{represents}, and
\texttt{diagnosed\_by}.  This permits retrieval by problem structure rather
than keyword overlap alone.

A procedural capsule is accepted only when it supplies all fields in
\begin{equation}
\mathcal{C} =
\{P,I,C,S,V,F,Q,T\},
\end{equation}
where $P$ is purpose, $I$ inputs, $C$ preconditions, $S$ executable steps,
$V$ invariants, $F$ failure signals, $Q$ verification properties, and $T$
transfer scope.  A capsule with any missing required field fails closed.

For example, the retained capsule for topological priority planning can be
summarised as
\begin{verbatim}
topological_priority_planning
  purpose: order dependent work safely
  inputs: task IDs, dependencies, priorities, capability levels
  preconditions: unique IDs; known dependency references
  invariants: dependency-before-dependent; exactly-once
  method: repeatedly select the highest-priority ready task
  failure signals: cycle; missing dependency; duplicate ID
  verification: ordering, gap insertion, and rejection properties
  transfer: build systems, project plans, workflow engines
\end{verbatim}

The original $5/6$ failure was preserved as an immutable negative receipt.  It
records the expected and observed interface behaviour without copying a valid
solution.  Thus the system can retain what failed and why without converting a
hidden examination into an answer book.

The promoted lexicon procedure was
\begin{verbatim}
procedure_functional_glyph_lexicon_v2
\end{verbatim}
and its initial cohort reconstructed both functional capsules, resolved all
provenance, rejected the deliberately corrupted capsule, and stored no raw
answer.

\subsection{Typed Method Compilation and Feedback Repair}

The second promoted component was
\begin{verbatim}
procedure_functional_glyph_method_compiler_v1
\end{verbatim}
which compiles verified procedural intermediate representations into
dependency-free Python.  The initial compiler supported 15 semantic operators:
eight for dependency-safe priority planning and seven for restartable mission
execution.

The planning operators validate unique identifiers and dependency references,
derive the dependency-ready set, order ready work by priority, insert a
learning action below the capability threshold, emit task identifiers, advance
the completed set, and reject a stalled cycle.  The runner operators return the
current task, reject an outcome for the wrong task, advance after success,
retry after failure, snapshot state, restore state, and detect completion.

Unknown, incomplete, or authority-expanding operators are rejected before
code generation.  Four of four corrupt or unsafe method programs were blocked,
and a renamed source-disjoint transfer contract passed.

The candidate-revision loop was also corrected.  The earlier bridge reduced
all hidden execution failures to the generic token
\texttt{executable\_contract\_failed}, discarding the bounded diagnostic
information needed for repair.  The revised bridge exposes only safe failure
metadata---return code, bounded output fragments, failed property names, and
counts---while continuing to withhold the hidden verifier implementation and
solution.  Candidate replay is bounded, passing components are preserved, and
the best failed candidate can seed the next private challenger.

\subsection{Unchanged Examination Recovery}

The original six-artifact examination was then rerun unchanged.  The
Functional Glyph Lexicon supplied the method structure, and the typed compiler
constructed the missing planner and runner without reopening the source lesson
or retrieving a raw program.  The final Python artifact returned
\texttt{SEALED\_MISSION\_EXECUTOR\_PASS} from the same independent contract.

\begin{table}[htbp]
\centering
\caption{Functional-memory reconstruction result.}
\label{tab:functional-memory-reconstruction}
\begin{tabular}{lr}
\hline
\textbf{Measure} & \textbf{Result} \\
\hline
Fresh competency artifacts & $6/6$ \\
Functional capsules & $2/2$ \\
Semantic operators compiled & $15$ \\
Source-disjoint method transfer & Pass \\
Corrupt or unsafe programs rejected & $4/4$ \\
Malicious mission candidates rejected & $4/4$ \\
Raw solutions retrieved & $0$ \\
Learning sources reopened & $0$ \\
Unsafe executions or live writes & $0$ \\
Focused regression tests & $13/13$ \\
\hline
\end{tabular}
\end{table}

The successful result promoted
\texttt{procedure\_open\_competency\_reconstruction\_mission\_v1}.  The
causal comparison was therefore direct:
\begin{equation}
\underbrace{\text{shallow topic glyphs}}_{5/6}
\quad\longrightarrow\quad
\underbrace{\text{typed functional capsules + method compiler}}_{6/6},
\end{equation}
under the unchanged hidden execution contract.

\subsection{Autonomous Functional-Memory Induction}

The initial repair demonstrated two hand-seeded procedural capsules, but scale
required a generic conversion from verified experience into functional memory.
The promoted procedure
\begin{verbatim}
procedure_autonomous_functional_glyph_induction_v1
\end{verbatim}
implements that conversion without a separate subject-specific loader.

For every candidate trajectory, the inducer requires:
\begin{enumerate}
    \item a complete source-closed capsule;
    \item resolved evidence anchors;
    \item successful fresh candidate execution;
    \item rejection of every supplied counterexample;
    \item rejection of every unsafe variant; and
    \item zero live-repository writes.
\end{enumerate}
Only then does it create a procedure node, grounded subskill concepts, a
verifier node, transfer edges, and provenance links.  Reapplying the same
trajectories is content-idempotent and creates neither duplicate nodes nor
duplicate edges.

Before graph insertion, the capsule and its fresh execution authority are
frozen together as a SHA-256-addressed immutable receipt:
\begin{equation}
r_i = H\!\left(\operatorname{canonical}
\left(C_i, E_i, Q_i, U_i\right)\right),
\end{equation}
where $C_i$ is the functional capsule, $E_i$ the independent execution
outcome, $Q_i$ the counterexample result, and $U_i$ the unsafe-variant result.
This prevents a later curriculum refresh from silently changing evidence
beneath an already retained glyph.  Six of six frozen receipts resolved to
their committed content.

\begin{table}[htbp]
\centering
\caption{Autonomous functional-memory induction.}
\label{tab:functional-memory-induction}
\begin{tabular}{lrr}
\hline
\textbf{Measure} & \textbf{Before} & \textbf{After} \\
\hline
Functional nodes & 11 & 83 \\
Typed relationships & 13 & 79 \\
Reconstructable procedures & 2 & 8 \\
Automatically induced procedures & 0 & 6 \\
Mixed-domain compositions passed & --- & $2/2$ \\
Corrupted capsules rejected & --- & $6/6$ \\
Raw answers stored & 0 & 0 \\
Original source documents opened & 0 & 0 \\
Unsafe actions or live writes & 0 & 0 \\
\hline
\end{tabular}
\end{table}

The induced procedures cover Algorithms and Data Structures, Software
Engineering, Mathematics, English natural-language cognition, Python, and
Scientific Method.  Two mixed-domain compositions reconstructed successfully:
\begin{align}
\text{algorithms}
&\rightarrow \text{software engineering}
\rightarrow \text{English communication}, \\
\text{mathematics}
&\rightarrow \text{scientific investigation}
\rightarrow \text{Python execution}.
\end{align}
Every component resolved through the typed graph with its original learning
sources closed.

\subsection{Memory-Qualified Advanced and Expert Evidence}

Competency level and reconstructability are deliberately separate authorities.
The progressive competency ledger alone may award a level; memory induction
cannot.  The additional memory-qualified status asks whether an already earned
competence can be reconstructed and applied after the original learning
material is unavailable.

At promotion time, the full ledger reported two Expert, six Advanced, and 11
Intermediate subjects.  Within the six-subject functional cohort, two Expert
and four Advanced subjects were additionally memory-qualified.

\begin{table}[htbp]
\centering
\caption{Source-closed memory qualification of existing competency levels.}
\label{tab:memory-qualified-levels}
\begin{tabular}{ll}
\hline
\textbf{Existing level} & \textbf{Memory-qualified subjects} \\
\hline
Expert & Algorithms and Data Structures; Software Engineering \\
Advanced & Mathematics; English; Python; Scientific Method \\
\hline
\end{tabular}
\end{table}

Five of the six functional-mastery reconstructions were source-disjoint.
Python passed functional reconstruction but remains explicitly marked as not
source-disjoint in that particular cohort.  No competency award was generated
by the induction process.

\subsection{Correction of the Advanced Python Timing Authority}

The closing regression exposed a separate false negative in the Advanced
Python executor.  A correct asynchronous pipeline was rejected when a busy
host exceeded a fixed 130\,ms wall-clock threshold, even though the probe
observed concurrent worker overlap, correct error propagation, and complete
cleanup.  The old assertion measured transient laptop load rather than the
semantic property of concurrent execution.

The repaired sealed authority uses observed worker overlap as the concurrency
criterion and retains a broad five-second bound only for deadlock or runaway
execution.  It separately verifies typing protocols, generators, asynchronous
context management, packaging, error propagation, cleanup, transfer, and
security.

\begin{table}[htbp]
\centering
\caption{Corrected Advanced Python execution authority.}
\label{tab:advanced-python-authority}
\begin{tabular}{lr}
\hline
\textbf{Measure} & \textbf{Result} \\
\hline
Advanced Python competency properties & $7/7$ \\
Source-disjoint transfer & Pass \\
Malicious variants rejected & $6/6$ \\
Final score & $1.0$ \\
Unsafe programs executed & $0$ \\
Live repository writes & $0$ \\
\hline
\end{tabular}
\end{table}

This correction changed the measurement authority, not the competency
thresholds and not the ledger evidence.

\subsection{Integrated Verification and Claim Boundary}

The complete closing regression passed 35 of 35 focused tests while the live
mastery service remained active with autonomous functional-memory induction
loaded.  The architecture now supports the loop
\begin{equation}
\text{verified experience}
\rightarrow \text{typed functional capsule}
\rightarrow \text{immutable evidence receipt}
\rightarrow \text{source-closed reconstruction}
\rightarrow \text{fresh execution and criticism}
\rightarrow \text{retained transferable method}.
\end{equation}

The result closes the immediate representation defect exposed by the original
$5/6$ examination.  It does not establish complete Python mastery, arbitrary
knowledge induction, or unrestricted self-programming.  The current induced
cohort contains six competency procedures and relies on installed bounded
execution authorities.  An unfamiliar subject must still acquire or invent a
safe executor and face an independent consequence before its trajectory can be
promoted into functional memory.

The remaining memory task is continual population from naturally arriving
verified trajectories across the wider curriculum.  The remaining progression
from Advanced to Expert is predominantly repeated unfamiliar practical work,
debugging recovery, source-disjoint transfer, low-scaffolding consistency, and
elapsed closed-book retention rather than another fundamental memory-format
redesign.

